%===========================================================
%== Präambel ===============================================
%===========================================================
\documentclass[
    paper=a4,
    bibtotocnumbered,
    liststotocnumbered,
    oneside,
    12pt,
    listof=totoc,
    toc=chapterentrywithdots,
    listof=entryprefix,
]{scrreprt}

\usepackage[a4paper, left=2.5cm, right=2.5cm, top=1.25cm, bottom=1.25cm, includehead, includefoot]{geometry}


\usepackage{float}
\usepackage[T1]{fontenc}%Encoding
\usepackage[german]{babel}%German-specific commmands
%\usepackage[utf8]{inputenc} %Wird nicht benötigt bei XeLatex
\usepackage{fancyhdr}%% Kopfzeile u. Fußzeile
\usepackage{fontspec}%% Schriftart
\usepackage[onehalfspacing]{setspace}%-- Zeilenabstand 1,5
\usepackage[toc, nopostdot, acronyms]{glossaries}%% Glossar
\usepackage[intoc]{nomencl}%% Symbolverzeichnis
\usepackage{etoolbox}%Sorgt dafür, das Symbolverzeichnis im Header steht
%\usepackage{listings}
\usepackage[]{minted}
\usepackage[titles]{tocloft}
\usepackage{scrhack}
\usepackage{colortbl}%Einfärben von Tabellen-Zellen, Zeilen, Spalten
\usepackage[]{layout}
\usepackage[]{blindtext}
%\usepackage{showframe}% zum Anzeigen des Seitenlayouts
\usepackage[figurename={Abb.},tablename={Tab.}]{caption}%Anpassung der Namen für Abbildung und Tabelle
\usepackage[]{url}
\usepackage[]{hyperref}
\usepackage{chngcntr}%% Nummerierung der Abbildungen und Tabellen fortlaufend %%%%%%
\usepackage{tocbibind}

% Paragraph styles
\setlength{\parindent}{0cm}
\setlength{\parskip}{6pt}

%% Kopfzeile u. Fußzeile
\pagestyle{fancy}
\fancyfoot{}
\fancyfoot[R]{\thepage}
\fancyhead{}
\fancyhead[L]{\nouppercase{\leftmark}}
\fancypagestyle{plain}{\pagestyle{fancy}}


%-- Standardschriftart auf Times New Roman umstellen
\setmainfont[ Path = Fonts/,
BoldFont=timesbd.ttf,
ItalicFont=timesi.ttf,
BoldItalicFont=timesbi.ttf
]{times.ttf}

%-- Kapitelüberschriften auf Standardschriftart(Times New Roman) umstellen
\setkomafont{disposition}{\normalcolor\bfseries}

% Entfernt "Kapitel X" aus der Kopfzeile vor der Kapitelüberschrift
\renewcommand{\chaptermark}[1]{\markboth{\MakeUppercase{#1}}{}}
 
%Abstände nach Kapiteln 
\RedeclareSectionCommand[%
  beforeskip=0pt,
  afterskip=1\baselineskip plus .1\baselineskip minus .167\baselineskip
]{chapter}

%% Nummerierung der Abbildungen und Tabellen fortlaufend %%%%%%
\counterwithout{figure}{chapter}
\counterwithout{table}{chapter}

%% Rahmen in Textbreite um Bilder
\floatstyle{boxed}
\restylefloat{figure}
\restylefloat{table}
\restylefloat{listing}

%% Captions linksbündig orientieren
\captionsetup{justification=raggedright,singlelinecheck=false}

% Einbinden von Glossar u. Abkürzungsverzeichnis
%-----------------------------------------------------------------------------
% Glossar
%-----------------------------------------------------------------------------

% \newglossaryentry{Literal}
% {
%     name=Literal,
%     description={\todo{Was ist ein Literal}}
% }
\input{Verzeichnisse/Abkürzungsverzeichnis.tex}
\makeglossaries

\begin{document}

%===========================================================
%== Titelseite =============================================
%===========================================================
\newgeometry{left=2.5cm, right=2.5cm, top=2.5cm, bottom=2.5cm}
\begin{titlepage}
    \begin{center}
    \includegraphics[width=\textwidth]{Images/logo}
        \begin{Large}
        Hochschule für angewandte Wissenschaften Coburg
        \\
        Fakultät Elektrotechnik und Informatik
        \par
        \end{Large}
        \vspace{2.0cm}
        
        \begin{Large}
            Studiengang: Informatik
            \par
        \end{Large}
        \vspace{1.5cm}
        
        {\Large
            Bachelorarbeit
        }
        \vspace{2.0cm}

        \begin{huge}
            \textbf{Vergleich von Google Go und Apple Swift} 
            \par
        \end{huge}        
        \vfill
        
        \begin{huge}
            Daniel Müller
        \end{huge}
        \vspace{2.0cm}
        
        \begin{large}
            Abgabe der Arbeit: 08. Juni 2017
            
            Betreut durch:
            \\
            Prof. Volkhard Pfeiffer, Hochschule Coburg
            \\
            %<optional: Zweitgutachter: Prof. Dr. XXX, Hochschule Coburg>
            \par
        \end{large}
        
    \end{center}
\end{titlepage}
\restoregeometry
    
%===========================================================
%== Inhaltsverzeichnis =====================================
%===========================================================   
\begin{singlespace}
\renewcommand{\cftchapleader}{\cftdotfill{\cftdotsep}} % Punkte im Inhaltsverzeichnis bei Kapiteln
\tableofcontents
\end{singlespace}

%===========================================================
%== Abbildungsverzeichnis ==================================
%===========================================================
\newpage
\renewcommand{\cftfigpresnum}{Abb. }
\renewcommand{\cftfigaftersnum}{:}
\setlength{\cftfignumwidth}{2cm}
\setlength{\cftfigindent}{0cm}
\listoffigures

%===========================================================
%== Tabellenverzeichnis ====================================
%===========================================================
\newpage
\renewcommand{\cfttabpresnum}{Tab. }
\renewcommand{\cfttabaftersnum}{:}
\setlength{\cfttabnumwidth}{2cm}
\setlength{\cfttabindent}{0cm}
\listoftables

%===========================================================
%== Codeverzeichnis ========================================
%===========================================================
\newpage
\renewcommand\listingscaption{Code}
\renewcommand\listoflistingscaption{Codebeispielverzeichnis}
\renewcommand{\cftfigpresnum}{Code }
\listoflistings
\addcontentsline{toc}{chapter}{Codebeispielverzeichnis}

%===========================================================
%== Symbolverzeichnis ======================================
%===========================================================
\newpage
%% Symbolverzeichnis
\makenomenclature
% This will add the units
%----------------------------------------------
\newcommand{\nomunit}[1]{%
\renewcommand{\nomentryend}{\hspace*{\fill}#1}}
%----------------------------------------------
\renewcommand{\nomname}{Symbolverzeichnis}
%Sorgt dafür, das Symbolverzeichnis im Header steht
\patchcmd{\thenomenclature}
  {\chapter*{\nomname}}% usually only \chapter*{\nomname} is issued
  {\chapter*{\nomname}\markboth{\MakeUppercase\nomname}{\MakeUppercase\nomname}}
  {}{}
%-----------------------------------------------------------------------------------
% Symbolverzeichnis
%-----------------------------------------------------------------------------------
% Bsp: \nomenclature[prefix]{symbol}{description}

\nomenclature[]{$c$}{Speed of light in a vacuum inertial system \nomunit{$299,792,458\, m/s$}}

\nomenclature[]{$i$}{Ganzzahlige Laufvariable \nomunit{Integer}}
\printnomenclature


%===========================================================
%== Abkürzungsverzeichnis ==================================
%===========================================================
\newpage
\printglossary[type=\acronymtype, title=Abkürzungsverzeichnis, toctitle=Abkürzungsverzeichnis]

%<<<<<<<<<<<<<<<<<<<<<<<<<<<<<<<<<<<<<<<<<<<<<<<<<<<<<<<<<<<<<<<<<<<<<<<<<<<<<<<<<<<<
% Text Einleitung, Hauptteil usw. 
%<<<<<<<<<<<<<<<<<<<<<<<<<<<<<<<<<<<<<<<<<<<<<<<<<<<<<<<<<<<<<<<<<<<<<<<<<<<<<<<<<<<<

\chapter{Einleitung}
\blindtext

\section{Vorstellung des Kontexts der Arbeit}
\blindtext

\section{Zielsetzung/Motivation}
\Blindtext

\section{Aufgabenstellung}
\blindtext

\section{Aufbau der Arbeit}
\blindtext

\blinddocument

\chapter{Beispiele}
\section{Code-Test}
\blindtext

    \begin{listing}[]
    \caption{C\# Hello World}
    \label{listing:helloWorld}
    \begin{minted}[]{csharp}
    // Hello1.cs
    public class Hello1
    {
       public static void Main()
       {
          System.Console.WriteLine("Hello, World!");
       }
    }
    \end{minted}
    \end{listing}

\blindtext
\newpage

\section{Glossar-Test}
Um einen Eintrag im Glossar zu erzeugen, muss der defninierte Glossar-Begriff z.B. \Gls{computer} aufgerufen werden

%Um eine Abkürzung zu verwenden wird \acrlong{gcd} benutzt.
Um eine Abkürzung zu verwenden wird \gls{gcd} benutzt.
\newpage


\section{Bild-Test}
\blindtext

	\begin{figure}[h]
	\centering
	\includegraphics[width=0.5\textwidth]{Images/latex}
	\caption{Ein Beispielbild}%
	\label{figure:latex}%
	\end{figure}

Eingebunden.
\blindtext
\newpage

\section{Tabelle-Test}
Hier kommt eine Tabelle.
Die Tabelle \ref{table:test} wird so referenziert.
\blindtext
 
\begin{table}[h]
\centering
\begin{tabular}{|c|c|c|c|} 
 \hline
 \rowcolor[gray]{0.75} \textbf{Col1} & \textbf{Col2} & \textbf{Col2} & \textbf{Col3} \\
 \hline
 1 & 6 & 87837 & 787 \\ 
 \hline 
 2 & 7 & 78 & 5415 \\
 \hline 
 3 & 545 & 778 & 7507 \\
 \hline
 4 & 545 & 18744 & 7560 \\
 \hline
 5 & 88 & 788 & 6344 \\
 \hline
\end{tabular}
\caption{Table to test captions and labels}
\label{table:test}
\end{table}

\blindtext

\newpage

\section{Literatur-Test}
kdksdlf \cite{brooks2001-silver} sklfdjslfdsl 

\chapter{Zusammenfassung}

\chapter{Ausblick}


%===========================================================
%== Glossar ================================================
%===========================================================
%% Glossar
\renewcommand*{\glossaryentrynumbers}[1]{} %Entfernt die Seitenzahl am Ende der Glossar-Beschreibung

\printglossary[title=Glossar,toctitle=Glossar]

%===========================================================
%== Literaturverzeichnis ===================================
%===========================================================
\bibliography{Literatur/literatur}
\bibliographystyle{alphadin}

%===========================================================
%== Anhang =================================================
%===========================================================
\appendix
\chapter{Anhang}
\newpage
\section{Ehrenwörtliche Erklärung}
\begin{centering}
\textbf{{\huge Ehrenwörtliche Erklärung}}
\par
\end{centering}

\vspace{2cm}

Ich versichere hiermit, dass ich meine/n Praxisbericht/Bachelorarbeit/Masterarbeit mit dem Titel

\vspace{2cm}

\begin{tabular*}{\linewidth}{@{\extracolsep{\fill}}ccc}
 \\ \hline
 \vspace{2cm}
 \\ \hline
\end{tabular*}

\vspace{2cm}

selbständig verfasst, keine anderen als die angegebenen Quellen und Hilfsmittel benutzt sowie nicht an anderer Stelle als Prüfungsarbeit vorgelegt habe.

\vfill

\begin{tabular}{ccc}
\cline{1-1}
\parbox{7cm}{\raggedright Ort} &
\parbox{3cm}{\raggedright} &
\parbox{7cm}{\raggedright} \\
\vspace{2.5cm} \\
\cline{1-1} \cline{3-3}
\parbox{7cm}{\raggedright Datum} &
\parbox{3cm}{\raggedright} &
\parbox{7cm}{\raggedright Unterschrift} \\ 
\end{tabular}

\end{document}
